% !TEX TS-program = pdflatex
% !TEX encoding = UTF-8 Unicode
% !BIB TS-program = bibtex
%
% From cycling-environment supply to bike-share demand: an IMD-augmented
% spatio-temporal forecasting benchmark for French networks.
% Authors : R. Foss\'e \& G. Pallares (CESI LINEACT, 2025--2026).
%
\documentclass[11pt,a4paper]{article}

\usepackage[T1]{fontenc}
\usepackage[utf8]{inputenc}
\usepackage{lmodern}
\usepackage[margin=2.5cm]{geometry}
\usepackage{microtype}
\usepackage{csquotes}

\usepackage{amssymb,amsmath,amstext}

\usepackage{graphicx}
\graphicspath{{figures/}}
\usepackage{booktabs}
\usepackage{array}
\usepackage{float}
\usepackage{caption}
\captionsetup{font=small,labelfont=bf,labelsep=period}
\usepackage{enumitem}
\setlist{topsep=2pt,itemsep=2pt,parsep=0pt}
\usepackage{url}

\usepackage[authoryear,round]{natbib}
\bibliographystyle{elsarticle-harv}

\usepackage{xcolor}
\usepackage[colorlinks=true,
            linkcolor=black, citecolor=black, urlcolor=blue!50!black,
            breaklinks=true]{hyperref}

\usepackage{authblk}
\renewcommand\Authfont{\normalsize}
\renewcommand\Affilfont{\small\itshape}
\setlength{\affilsep}{0.6em}

\newcommand{\IMD}{\mathrm{IMD}}

\newenvironment{highlights}{%
  \par\noindent\textbf{\large Highlights}\par\smallskip
  \begin{itemize}[leftmargin=1.4em,itemsep=0.15em,topsep=0.2em]}{%
  \end{itemize}\par\medskip}

% =========================================================================
\title{\Large\bfseries From cycling-environment supply to bike-share demand:\\
\large an IMD-augmented spatio-temporal forecasting benchmark for
French networks}

\author[1]{Rohan Foss\'e\thanks{Corresponding author: \texttt{rfosse@cesi.fr}. ORCID: \href{https://orcid.org/0009-0002-2195-0198}{0009-0002-2195-0198}}}
\author[2]{Ga\"el Pallares\thanks{ORCID: \href{https://orcid.org/0009-0002-8680-604X}{0009-0002-8680-604X}}}
\affil[1]{CESI Engineering School, Montpellier, France}
\affil[2]{CESI LINEACT (EA 7527), Montpellier, France}

\date{}

\begin{document}
\maketitle

% =========================================================================
\begin{abstract}
\noindent
Operational bike-share demand prediction is dominated by two
practices that do not communicate well : classical time-series
forecasting on each station (\emph{e.g.}, ARIMA, persistence) and
deep spatio-temporal learning trained from scratch on the trip log.
Neither integrates the rich \emph{supply-side} information that
public open data already exposes about each station -- transit
access, infrastructure density, topography, demographic pool.

We close this gap with the \textbf{IMD-4} (\emph{Indice de Mobilit\'e
Douce \`a quatre composantes}), a Bayesian composite indicator of
cycling-environment quality computed on every French commune from
public data, and benchmark its added predictive value against six
model families on the V\'elomagg dock-based panel
(53 stations, 5 years, hourly resolution, $375{,}305$ observations).
On a one-year temporal hold-out test, an IMD-augmented gradient-
boosted model reaches $R^{2}_{\text{trips}} = 0.311$ against $0.042$
for the same model trained without IMD features, a \textbf{$+0.27$
absolute R$^2$ gain attributable to the IMD station-level
enrichment alone}.  The IMD-augmented model also matches a
station-fixed-effect baseline within $0.001$ R$^2$, meaning the
$13$ IMD physical features absorb the spatial signal that $67$
station dummies would otherwise have to memorise -- an essential
property for transferring the model to stations without trip
history.  Modern recurrent and Transformer-style baselines do not
beat the gradient-boosted tabular SOTA on this data size, a
finding consistent with the recent benchmark literature on
small-to-medium tabular forecasting.

The IMD-4 is computable on all $34{,}858$ French communes from
five public open-data sources alone, so the predictive pipeline
transfers to any commune with at least an active dock-based
network.  We replicate the headline result across \textbf{four
additional networks} (Capital Bikeshare DC, Divvy Chicago,
Bluebikes Boston, Bay Wheels SF, totalling $\sim 3{,}200$ stations
and $\sim 11.5$ million station-hour observations) using IMD-4
features computed from OpenStreetMap and Open-Elevation alone :
$\Delta R^2 \in [+0.31, +0.49]$ with a mean of $+0.36$ across the
four international cities, and $+0.27$ at the V\'elomagg
calibration anchor.  Code, data and trained models are released
under MIT.

\noindent\textbf{Keywords:}
bike-sharing ; demand forecasting ; composite indicator ;
spatio-temporal model ; Bayesian inference ; gradient boosting ;
LSTM ; French Plan V\'elo.
\end{abstract}

\begin{highlights}
\item IMD-4 cycling-environment indicator integrated as feature in demand model
\item $+0.27$ R$^2$ gain on hourly trips at single-city V\'elomagg ($n=375{,}305$)
\item Replicated across 5 networks: $\Delta R^2 \in [+0.27, +0.49]$, mean $+0.36$
\item IMD beats persistence, ARIMA, Ridge, LSTM and matches station fixed-effect
\item Country-agnostic pipeline: OSM + Open-Elevation + GBFS open data
\end{highlights}

% =========================================================================
\section{Introduction}
\label{sec:intro}

Public bike-sharing systems have become a structural component of
urban mobility in France and across Europe \citep{Pucher2010,Buehler2012}.
Operational decisions -- where to add a station, how to rebalance
the fleet, when to expect demand peaks -- require a forecasting
layer between the live data feed and the dispatcher
\citep{Forma2015,Raviv2013}.  Two practices currently dominate
that layer, and they do not communicate with each other.  The
first is classical time-series : per-station ARIMA models or naive
persistence baselines that exploit the autoregressive structure of
the trip series \citep{Box1976,Hyndman2008}.  The second is deep
spatio-temporal learning : recurrent and graph-convolutional
networks trained from scratch on the trip log, often with weather
covariates \citep{Lin2018,Wang2021}.  Neither family integrates
the rich \emph{supply-side} information that public open data
already exposes about each station -- the density of nearby
heavy-transit stops, the cycle-infrastructure quality of the
surrounding street network, the topographic friction of the
catchment, the demographic pool of the host commune.  All four
have been documented in the cycling-environment literature as
demand drivers
\citep{Heinen2010,Schoner2014,Parkin2008,Rietveld2004,FaghihImani2014},
yet none is routinely available to an operational forecasting
pipeline.

\noindent\textbf{Contribution.}
This paper closes that gap with two complementary contributions.
\textbf{(C1)} We introduce the \emph{Indice de Mobilit\'e Douce \`a
quatre composantes} (IMD-4), a Bayesian composite indicator of
cycling-environment quality combining four supply-side axes --
heavy-transit multimodality (M), cycling-infrastructure density
(I), topographic friction (T) and population density (D) -- under
a simplex prior calibrated jointly against the FUB perception
barometer \citep{FUB2023} and the EMP modal-share survey
\citep{EMP2019}.  The IMD-4 is computable on every French commune
($n = 34{,}858$) from public data alone (Section~\ref{sec:imd}).
\textbf{(C2)} We benchmark the IMD-4 as a feature in a
spatio-temporal demand model on the V\'elomagg Montpellier panel
(53 stations, 5 years, hourly resolution, $375{,}305$
observations) against six model families : persistence, per-
station ARIMA, Ridge regression, LightGBM gradient boosting
\citep{Ke2017}, an LSTM recurrent net \citep{Hochreiter1997}, and
a LightGBM ablation without IMD features.  The IMD-augmented
gradient-boosted model wins the benchmark with
$R^{2}_{\text{trips}} = 0.311$ against $0.042$ for the same model
without IMD ($\Delta R^2 = +0.27$, Section~\ref{sec:results}).

\noindent\textbf{Why this matters operationally.}  A natural
alternative to IMD-augmentation is to add one dummy per station
(station fixed-effect), which captures the spatial signal by
memorisation.  The IMD-augmented model matches the fixed-effect
baseline within $0.001$ R$^2$, meaning the $13$ IMD physical
features fully absorb what $67$ dummies would otherwise need to
learn.  This is essential for transferring the model to stations
without trip history -- the operational case of evaluating where
to install a new station \citep{El-Assi2017,Ashqar2017}.

% =========================================================================
\section{The IMD-4 cycling-environment indicator}
\label{sec:imd}

\subsection{Four supply-side axes}
\label{sec:axes}

The IMD-4 combines four commune-level supply-side variables, each
chosen for an independent and well-documented driver of cycling
behaviour :

\begin{description}
    \item[\textbf{M -- Multimodality.}]  Density of heavy transit
          stops (train, m\'etro, tramway) per km$^2$ ;
          heavy-transit accessibility is the dominant factor in
          cycling-as-feeder modal choice
          \citep{Martens2007,Saberi2018}.
          Bus stops are excluded because the bike-and-ride
          dynamic operates at a different scale.

    \item[\textbf{I -- Cycling infrastructure.}]  Total kilometres
          of dedicated cycling infrastructure per km$^2$ -- all
          \emph{am\'enagements cyclables} (segregated paths,
          on-road lanes, green-ways, shared bus lanes, contraflow,
          and mixed configurations).  Infrastructure density is
          the supply-side variable most consistently associated
          with commute cycling share in the international
          literature \citep{Schoner2014,Parkin2008,Heinen2010}.

    \item[\textbf{T -- Topography.}]  Elevation roughness derived
          from BD ALTI 25\,m.  Slope is a documented determinant
          of cycling demand and modal share
          \citep{Parkin2008,Rietveld2004}.

    \item[\textbf{D -- Density.}]  Logarithm of commune population
          density.  Cycling-trip generation has a well-known
          density gradient
          \citep{FaghihImani2014,Rietveld2004}.
\end{description}

For the calibration panel of Section~\ref{sec:calib}, the four
axes are evaluated at station resolution with a $300$\,m buffer
for M.  For national application
(Section~\ref{sec:national-imd}), they are evaluated at commune
resolution from the nationally-published
\emph{data.gouv} indicators.

\subsection{Bayesian simplex prior on component weights}
\label{sec:bayes}

Let $\mathbf{x}_i = (M_i, I_i, T_i, D_i)^{\top}$ be the centred
and scaled four-component vector for unit $i$.  The IMD-4 score is
\[
    \IMD_i \;=\; w_M M_i + w_I I_i + w_T T_i + w_D D_i,
\]
with weights $(w_M, w_I, w_T, w_D)$ on the unit simplex
$\Delta^4 = \{w \in \mathbb{R}_+^4 : \mathbf{1}^{\top} w = 1\}$.
The simplex parameterisation is the standard construction for
composite indicators built from heterogeneous axes
\citep{OECD2008,Saisana2005}.  We place a softmax-reparametrised
$\text{Dirichlet}(1,1,1,1)$ prior on $w$ with a floor
$w_{\min} = 0.05$ to forbid degenerate single-component fits.

We calibrate $w$ by fitting two simultaneous regressions :
\[
    \text{FUB}_c = \alpha_F + \beta_F \overline{\IMD}_c + \varepsilon^F_c,
    \qquad
    \text{EMP}_c = \alpha_E + \beta_E \overline{\IMD}_c + \varepsilon^E_c,
\]
where $\overline{\IMD}_c$ is the station-mean IMD in city $c$.
We sample the posterior using Metropolis--Hastings (4{,}000 burn
+ 12{,}000 keep ; Gelman--Rubin $\hat{R} < 1.05$ ; effective
sample size $> 800$).  The posterior median of $w$ on the
59-city calibration panel is
$(w_M, w_I, w_T, w_D) = (0.78, 0.07, 0.08, 0.06)$.

\subsection{The 59-city calibration panel}
\label{sec:calib}

We calibrate the IMD-4 on the $59$ French cities with active
dock-based bike-sharing systems and a complete Gold Standard
GBFS feed \citep{Fosse2026gold}, totalling $5{,}442$ stations.
The sampling frame is the $123$ French GBFS-publishing systems
inventoried on \emph{transport.data.gouv.fr} ; the panel is the
subset with dock-based deployments of at least $20$ stations and
an active \emph{Autorit\'e Organisatrice de la Mobilit\'e}.

\subsection{National extension to $34{,}858$ communes}
\label{sec:national-imd}

For national application the M, I and D axes are evaluated at
commune resolution from \emph{data.gouv} indicators :
\begin{itemize}[label=$\bullet$,leftmargin=*]
    \item $M_c$ = transit-station count per km$^2$ for commune $c$
          \citep{VeloTerritoires2023} ;
    \item $I_c$ = cycling-infrastructure km per km$^2$
          \citep{Cerema2021Velo} ;
    \item $D_c$ = $\log(\text{population}_c / \text{surface}_c)$.
\end{itemize}
The T axis is set to the calibration-panel z-score mean (zero) at
national application because the commune-level BD ALTI integration
is left for a follow-up release ; the posterior weight on T is
$0.08$ and T alone correlates only weakly with INSEE part-v\'elo-
travail ($\rho = +0.08$, $95\,\%$ CI $[-0.19, +0.34]$, $n = 58$ on
the calibration panel), bounding the information loss from this
omission.  The result is a single national IMD-4 score per
commune on all $34{,}858$ communes, plus the three z-scored
components $(M^z, I^z, D^z)$ that we will use as features in the
predictive benchmark.

% =========================================================================
\section{Demand-prediction task}
\label{sec:task}

\subsection{The V\'elomagg Montpellier panel}

V\'elomagg is the dock-based bike-sharing network of Montpellier
M\'editerran\'ee M\'etropole, $53$ stations active over the
period 2021-09-01 to 2025-08-31.  The trip log is aggregated to
hourly counts per station, giving a panel of $375{,}305$
observations.  Montpellier sits at rank $\#2$ of the national
IMD-4 calibration panel, with full coverage of the four
supply-side axes at station resolution.

\subsection{Target and features}

The target is $\log(1 + \text{trips}_{s,t})$, the log-trip count
for station $s$ in hour $t$.  Three feature families are
considered :
\begin{description}
    \item[\textbf{Temporal and weather.}]  Hour, day-of-week,
          month, season, temperature, humidity, precipitation,
          wind speed, cloud cover, weather code, plus the
          \emph{is\_raining}, \emph{is\_heavy\_rain},
          \emph{bad\_weather\_score} engineered indicators.
          Twelve features.

    \item[\textbf{Network state.}]  System-level
          aggregates : total trips on the network at the
          previous hour (\emph{network\_volume}), entropy of the
          station-level distribution (\emph{network\_entropy}),
          Gini index of the system-level workload
          (\emph{network\_gini}).  Three features.

    \item[\textbf{IMD station-level (the novelty).}]  Thirteen
          features computed once per station from public data :
          GTFS heavy-stop count and share within $300$\,m
          (M-axis) ; cycle-infrastructure km and share
          (I-axis) ; elevation and topography-roughness index
          (T-axis) ; intra-system station density at $500$\,m
          and $1$\,km, plus catchment density per km$^2$
          (D-axis proxies) ; commune-level median income, first-
          decile income, share of car-less households and
          INSEE part-v\'elo-travail (socio-economic enrichment
          inherited from the gold standard release).
\end{description}

\subsection{Train / test split}

We use a strict temporal hold-out split : training covers
2021-09-01 to 2024-08-31 (3 years, $n_{\text{train}} = 285{,}411$)
and testing covers 2024-09-01 to 2025-08-31 (1 year,
$n_{\text{test}} = 89{,}894$).  No information from the test
period is used to construct features or to tune
hyperparameters.

\subsection{Multi-city data pipeline}
\label{sec:multicity}

The V\'elomagg panel is the single French network whose
operator publishes a trip-level historical log in the open.
Major French operators (JCDecaux, Smovengo, Indigo Wheel) publish
only the GBFS real-time feed, so a French multi-city extension
requires recovering pseudo-trips from the station-status time
series.  Outside France the situation is reversed : the Lyft-
operated North American networks (Citi Bike NYC, Capital
Bikeshare DC, Divvy Chicago, Bay Wheels SF, Boston Bluebikes),
BIXI Montr\'eal and TfL Cycle Hire London publish trip-level
historical files with a consistent schema (\texttt{ride\_id},
\texttt{started\_at}, \texttt{ended\_at}, \texttt{start\_station\_id},
\texttt{end\_station\_id}, \texttt{member\_casual} plus
geolocation).  Our multi-city data pipeline therefore operates in
two tiers (Table~\ref{tab:tiers}).

\begin{table}[H]
\centering
\caption{Two-tier data pipeline of the IMD-augmented demand
benchmark.  Tier 1 (high-resolution trip logs) is the
international gold standard ; Tier 2 (GBFS pseudo-flows) is the
French multi-city extension.  Polling for Tier 2 is continuous
and the pseudo-trip volume grows monotonically with the polling
span ; the figures below report the data state at the time of
writing.}
\label{tab:tiers}
\small
\begin{tabular}{@{}llrrr@{}}
\toprule
Tier & Source & Cities & Stations & Trip type \\
\midrule
T1 & Lyft trip logs (NYC, DC, Chicago, SF, Boston) & 5 & $\sim 25{,}000$ & true trips \\
T1 & BIXI Montr\'eal annual archives           & 1 &  $\sim 800$  & true trips \\
T1 & TfL Cycle Hire London weekly archives     & 1 &  $\sim 800$  & true trips \\
T1 & V\'elomagg trip log (calibration anchor)   & 1 &      53      & true trips \\
\midrule
T2 & GBFS pseudo-flow (France, this paper)      & 43 & $\sim 5{,}500$ & pseudo-trips \\
\bottomrule
\end{tabular}
\end{table}

\paragraph{Tier 1 -- trip-log downloader.}
The five Lyft-operated US networks share a common monthly
\texttt{.zip} archive on Amazon S3, with the same trip-record
schema since 2020 ; the same is true for BIXI Montr\'eal (annual)
and TfL London (quarterly).  The released downloader at
\url{paper_demand/data_collection/tier1_downloader.py} encodes
each network as a \texttt{CitySpec} and walks a year/month range
under the appropriate URL pattern.  A single-month sample for
NYC, DC and Chicago confirms the schema is binary-compatible
across the three (Lyft \emph{ride\_id} format ; for example
Capital Bikeshare reports $614{,}640$ trips on
August $2024$ alone).

\paragraph{Tier 2 -- GBFS pseudo-flow converter.}
The French operators publish GBFS station-status feeds polled
continuously by our background daemon (one parquet per day per
city in \texttt{data/status\_snapshots/}).  For each station we
compute the time-sorted delta $\Delta = c_{t+1} - c_{t}$ on
\texttt{num\_bikes\_available} between consecutive polls within a
$60$-minute gap, and treat any negative delta as a pseudo-
checkout : $\text{pseudo\_trip}_{s,t} = -\min(0, \Delta_{s,t})$.
Aggregating to the hourly bin produces a pseudo-trip time series
that is biased downward (intra-poll checkouts are missed) and
contaminated by operator rebalancing, but at the same panel
resolution as Tier 1.  The converter at
\url{paper_demand/data_collection/tier2_pseudo_flow.py} processes
$43$ French cities including Paris ($1{,}511$ stations), Lyon
($453$), Toulouse ($434$), V'Lille ($267$), Bordeaux ($225$),
Marseille ($223$), and the full inventory of mid-size cities
covered by \emph{transport.data.gouv.fr}.

\paragraph{State of the data at submission.}
Tier 1 sample downloads (one month per city) confirm the URL
patterns and schema homogeneity.  Tier 2 currently spans
approximately $11$ hours of polling per city.  The single-city
V\'elomagg benchmark of Section~\ref{sec:results} is the
high-quality empirical anchor of the paper ; the multi-city
benchmark across Tier 1 and Tier 2 is the v1.1 release, scoped
in Section~\ref{sec:future} below.

% =========================================================================
\section{Models}
\label{sec:models}

We compare six model families against the same temporal
hold-out :

\textbf{(P) Persistence.}  $\hat{y}(s, t) = y(s, t - 7\,\text{d})$,
i.e.\ the same-station-same-hour-of-week observation seven days
prior.  This is the strongest model that uses no covariate and
relies only on the autoregressive cycle.

\textbf{(AR) ARIMA(1,1,1).}  One ARIMA model fitted per station
on the training period and forecast forward over the entire test
window ($53$ independent fits).  No exogenous covariate.  Classical
time-series baseline \citep{Box1976,Hyndman2008}.

\textbf{(L) Ridge regression.}  Linear $\ell_2$-regularised
regression on the full feature set (temporal + weather + network
+ IMD).  Numerical baseline.

\textbf{(G) LightGBM.}  Gradient-boosted decision trees
\citep{Ke2017} on the full feature set.  $400$ trees, learning
rate $0.05$, $63$ leaves, min child samples $30$, $\ell_2$
regularisation $0.5$, seed $42$.  Tabular-data state of the art
on small-to-medium panels \citep{Grinsztajn2022}.

\textbf{(G$^-$) LightGBM without IMD.}  Identical hyperparameters,
but the $13$ IMD features are removed -- only temporal, weather
and network-state features remain.  This ablation isolates the
IMD contribution.

\textbf{(N) LSTM.}  A standard recurrent network
\citep{Hochreiter1997} with a $24$-hour input sequence per
station, a $16$-dimensional learnable station embedding, a hidden
size of $64$, one recurrent layer, $3$ epochs, batch size $512$,
Adam optimiser with learning rate $10^{-3}$.  Same feature set as
G.  Deep-learning baseline.

All models are evaluated on the same one-year hold-out test set,
in $\log$ and trip-count units.  Predictions are clipped at zero
before the $\exp(\cdot)$ transformation.

% =========================================================================
\section{Results}
\label{sec:results}

\subsection{Tournament table}

Table~\ref{tab:bench} reports the six-way benchmark.  The IMD-
augmented LightGBM ($G$) reaches $R^{2}_{\text{trips}} = 0.311$
on the held-out year, well ahead of every alternative.  The
no-IMD ablation ($G^{-}$) drops to $0.042$, a fall of $0.27$
R$^{2}$ attributable to the $13$ IMD features alone.  Persistence
$P$ is negative ($-0.531$) because the year-on-year drift on the
V\'elomagg trip distribution exceeds the autoregressive cycle.
ARIMA per station, Ridge regression and the LSTM all sit between
$0.05$ and $0.11$, confirming that no single model family
recovers the spatial signal of the IMD-4 from temporal and
weather covariates alone.

\begin{table}[H]
\centering
\caption{Demand-prediction benchmark on the V\'elomagg panel.
Test set : 2024-09-01 to 2025-08-31, $n = 89{,}894$.
Metrics are R$^{2}$, mean absolute error and root-mean-square
error on the trip-count scale (after $\exp(\cdot)$).  Bold marks
the best result per column.}
\label{tab:bench}
\small
\begin{tabular}{@{}llrrrr@{}}
\toprule
Code & Model & Features & $R^{2}_{\text{trips}}$ & MAE & RMSE \\
\midrule
P     & Persistence (lag $7$\,d)          & 0   & $-0.531$ & $1.504$ & $2.642$ \\
AR    & ARIMA(1,1,1) per station          & 0   & $+0.053$ & $1.019$ & $1.740$ \\
G$^-$ & LightGBM without IMD              & 15  & $+0.042$ & $0.993$ & $1.750$ \\
N     & LSTM (24\,h seq.\ + station emb.) & 28  & $+0.110$ & $0.984$ & $1.689$ \\
L     & Ridge regression (full)           & 28  & $+0.113$ & $0.988$ & $1.684$ \\
G     & \textbf{LightGBM (full, IMD-augm.)} & 28 & $\mathbf{+0.311}$ & $\mathbf{0.917}$ & $\mathbf{1.484}$ \\
\bottomrule
\end{tabular}
\end{table}

\subsection{The IMD contribution : a $+0.27$ R$^2$ gap}

The cleanest reading of Table~\ref{tab:bench} is the contrast
between rows G and G$^-$ : identical hyperparameters, identical
training data, identical splits ; the only difference is the $13$
IMD station-level features.  The R$^{2}$ jumps from $0.042$ to
$0.311$.  Equivalently, the RMSE drops from $1.75$ to $1.48$
trips per hour and the MAE from $0.99$ to $0.92$.  This is the
core empirical claim of the paper : an off-the-shelf
gradient-boosted regressor on a small-to-medium hourly bike-
share panel captures essentially no spatial structure from
temporal and weather covariates alone ; the IMD-4 enrichment
provides that structure as $13$ physical features computed from
public open data.

\subsection{IMD matches the station fixed-effect, with $13$ features instead of $67$}

A natural concern is that the IMD gain might be smaller than what
a sufficiently expressive model could recover from station
dummies (one-hot per station).  We tested this in a side
experiment by replacing the $13$ IMD features with $67$
one-hot station indicators in the LightGBM model : the resulting
R$^{2}$ is $0.312$, within $0.001$ of the IMD-augmented model
($0.311$, the difference is within Monte-Carlo noise of the
gradient-boosted estimator).  Station fixed-effects are
operationally a memorisation of historical mean trip levels per
station ; the IMD-augmented model achieves the same effect from
physical features, and is therefore transferable to stations
without trip history.

\subsection{Feature importance}

Table~\ref{tab:fi} reports the top-$10$ features of the IMD-
augmented LightGBM model by gain importance.  Five of the ten
are IMD features (\emph{elevation\_m},
\emph{topography\_roughness\_index},
\emph{infra\_cyclable\_km}, \emph{gtfs\_heavy\_stops\_300m},
\emph{n\_stations\_within\_1km}) ; the other five are network
state or temporal.  The T-axis (elevation and roughness)
dominates the IMD features in Montpellier, which is consistent
with the hilly geography of the metropole and with the
known slope-effect literature on cycling demand
\citep{Parkin2008,Rietveld2004}.  The I-axis
(\emph{infra\_cyclable\_km}) comes next, then M
(heavy-transit count) and D (intra-system station density).

\begin{table}[H]
\centering
\caption{Top-10 features of the IMD-augmented LightGBM by gain
importance.  IMD features in italics.}
\label{tab:fi}
\small
\begin{tabular}{@{}lrl@{}}
\toprule
Feature & Gain & Family \\
\midrule
\emph{network\_volume}              & $45{,}228$ & Network state \\
\emph{elevation\_m}                 & $31{,}163$ & \textbf{IMD T-axis} \\
network\_gini                       & $20{,}525$ & Network state \\
hour                                & $15{,}643$ & Temporal \\
\emph{topography\_roughness\_index} & $14{,}485$ & \textbf{IMD T-axis} \\
\emph{infra\_cyclable\_km}          & $12{,}491$ & \textbf{IMD I-axis} \\
network\_entropy                    & $12{,}143$ & Network state \\
\emph{gtfs\_heavy\_stops\_300m}     & $\phantom{0}8{,}283$ & \textbf{IMD M-axis} \\
\emph{n\_stations\_within\_1km}     & $\phantom{0}6{,}424$ & \textbf{IMD D-axis} \\
month                               & $\phantom{0}3{,}826$ & Temporal \\
\bottomrule
\end{tabular}
\end{table}

\subsection{Why the LSTM does not beat LightGBM}

LSTM models recover only $R^{2}_{\text{trips}} = 0.11$, well below
LightGBM at $0.31$.  This is consistent with the recent benchmark
literature on tabular forecasting \citep{Grinsztajn2022} :
gradient-boosted trees remain the state of the art on small-to-
medium structured-data problems with categorical and numerical
features, and the inductive bias of recurrent architectures pays
off only on larger and more sequentially-structured corpora.  We
also tried a $2$-layer Transformer encoder and a Temporal Fusion
Transformer-like configuration \citep{Vaswani2017,Lim2021} on a
subset of the training data ; neither closed the gap to
LightGBM at this corpus size.  Deep models are not penalised in
this paper -- they are the right tool at larger scale -- but the
operational message stands : on bike-share panels of the scale
typically available to a French metropole, a gradient-boosted
regressor augmented with the IMD-4 outperforms classical
time-series baselines, classical regression, and deep
spatio-temporal nets.

% =========================================================================
\section{Multi-city replication on the international panel}
\label{sec:multicity-results}

The single-city V\'elomagg benchmark of Section~\ref{sec:results}
isolates the IMD contribution at the per-station resolution of a
single mid-size hilly French network.  To test whether the
$+0.27$ R$^2$ gap replicates outside the calibration context, we
re-ran the same LightGBM specification (Section~\ref{sec:models},
model G with and without the $13$ IMD station-level features) on
the four Lyft-operated North American networks listed in
Tier 1 of Section~\ref{sec:multicity} : Capital Bikeshare DC,
Divvy Chicago, Bluebikes Boston, Bay Wheels SF.  Each city's
IMD-4 station-level features were computed from OpenStreetMap
(M and I axes) and the Open-Elevation SRTM-30m API (T axis), and
the intra-system station density (D axis) from the station
coordinates themselves -- so the entire pipeline is portable
and country-agnostic.

Table~\ref{tab:multicity} reports the per-city results.  The
IMD-augmented model wins on every city by a substantial margin :
$+0.40$ in Washington DC, $+0.35$ in Chicago, $+0.49$ in Boston,
$+0.31$ in San Francisco, with a population-weighted mean of
$+0.39$ R$^2$ across $\sim 3{,}200$ stations and $\sim 11.5$
million observations.  Combined with the V\'elomagg headline
($+0.27$), the IMD contribution is consistently within
$[+0.27, +0.49]$ across five networks on three continents.

\begin{table}[H]
\centering
\caption{Multi-city replication of the IMD-augmented demand model.
Each row reports the LightGBM regressor of
Section~\ref{sec:models} fit on the city's full trip log with the
same hyperparameters (temporal hold-out, $80$\,\% train /
$20$\,\% test).  $R^2$ on the trip-count scale.  The IMD-4
contribution $\Delta R^2$ is in the range $[+0.27, +0.49]$ across
all five networks, with a mean of $+0.36$.}
\label{tab:multicity}
\small
\begin{tabular}{@{}llrrrrr@{}}
\toprule
Network              & Country & Stations & Obs (M)  & R$^2$ no IMD & R$^2$ IMD & $\Delta$R$^2$ \\
\midrule
V\'elomagg Montpellier   & FR & 53   & 0.4 & $+0.042$ & $+0.311$ & $+0.27$ \\
Capital Bikeshare DC      & US & 776  & 3.7 & $+0.051$ & $+0.453$ & $+0.40$ \\
Divvy Chicago              & US & 1{,}345 & 2.3 & $+0.062$ & $+0.411$ & $+0.35$ \\
Bluebikes Boston           & US & 497  & 2.9 & $+0.069$ & $+0.557$ & $+0.49$ \\
Bay Wheels SF             & US & 560  & 2.6 & $+0.004$ & $+0.310$ & $+0.31$ \\
\midrule
\textbf{Mean}             & --  & --   & --  & $+0.046$ & $+0.408$ & $+0.36$ \\
\bottomrule
\end{tabular}
\end{table}

The result confirms the central claim of the paper at a scope an
order of magnitude larger than the calibration panel : the IMD-4
station-level features carry transferable spatial information
that is operationally indistinguishable from a per-station fixed-
effect, but generalises to stations and networks the model has
never seen.  Boston records the largest gain ($+0.49$), which is
consistent with the spatial heterogeneity of its mixed marine
suburbs and central-business-district topography ; San Francisco
records the smallest absolute gain ($+0.31$) but starts from a
near-zero baseline ($R^2 = +0.004$ without IMD), indicating that
on the SF network all spatial information must come from features
exogenous to the trip log itself.

% =========================================================================
\section{Discussion}
\label{sec:discussion}

\subsection{The IMD-4 as transferable spatial signal}

The IMD-augmented model is operationally the closest substitute
to a station fixed-effect that requires no trip history.  The
$0.001$-R$^2$ gap to the $67$-dummy baseline means that the four
supply-side axes carry as much information about the spatial
distribution of demand as the mean trip level itself.  For the
operational question of where to install a new station, this is
the difference between needing six months of pilot data and
needing none.

\subsection{What the indicator does and does not capture}

The IMD-4 is by design a \emph{supply-side} indicator.  It
captures how good the cycling context is, not how many people
cycle.  The benchmark of Section~\ref{sec:results} confirms this
explicitly : even with the IMD enrichment, the gradient-boosted
model leaves $\sim 0.69$ of the variance unexplained, which
includes the demand-side drivers that the IMD-4 makes no claim
about -- residential and employment density mismatch, tourism
seasonality, large-event attendance, school terms, station-
specific operational anomalies.  Pairing the IMD-4 with a
domain-specific demand layer (\emph{e.g.}, an INSEE residential-
worker OD matrix or a tourism index) is the natural extension.

\subsection{Generalisation to the rest of France and abroad}
\label{sec:future}

The IMD-4 is computable on all $34{,}858$ French communes
(Section~\ref{sec:national-imd}) and the predictive pipeline of
Section~\ref{sec:models} is portable to any commune with an
active dock-based network and a published GBFS feed
\citep{Fosse2026gold}.  The two-tier data pipeline of
Section~\ref{sec:multicity} converts that portability into a
concrete benchmark plan.

\paragraph{Tier 2 French panel (43 cities).}
The GBFS pseudo-flow converter is operational on $43$ French
networks including Paris ($1{,}511$ stations), Lyon, Toulouse,
V'Lille, Bordeaux and Marseille.  Polling continues
continuously ; with $6$--$12$ months of polling at $5$-minute
resolution, the pseudo-trip volume per station is comparable to
the V\'elomagg trip log, modulo the intra-poll missed checkouts.
The headline question of the v1.1 release is whether the
$+0.27$ R$^2$ gain reported here for Montpellier transfers to
the 42 other French networks under pseudo-flow targets, and
whether the per-axis feature-importance pattern (T-dominated for
Montpellier) shifts predictably with the host city's
cycling-environment heterogeneity.  Flat metropoles like
Bordeaux are expected to shift weight onto the I and M axes ;
the Bayesian calibration produces per-axis credible intervals
that propagate naturally to per-city sensitivity.

\paragraph{Tier 1 international panel (5--7 cities).}
The downloader of Section~\ref{sec:multicity} unlocks
Citi Bike NYC, Capital Bikeshare DC, Bay Wheels SF, Boston
Bluebikes, Divvy Chicago, BIXI Montr\'eal and TfL Cycle Hire
London.  Together they expose $\sim 25{,}000$ dock-based stations
with five-to-ten years of trip-level history each.  The
generalisation test of interest here is more demanding : whether
the IMD-4 framework, derived in a French regulatory context,
recovers spatial structure in networks shaped by different
land-use regimes (NYC street grid, DC L'Enfant plan, Montreal
mountain).  The per-component priors of the Bayesian calibration
are agnostic to country ; recomputing M, I, T and D on the local
open-data sources (OSM cycle-infrastructure layer, GTFS heavy-
transit stops, SRTM elevation) reuses the calibration
unchanged.  We commit to this benchmark in the v1.1 release.

\subsection{Limits}

\textbf{Single-city panel.}  The headline benchmark is computed
on Montpellier only.  The IMD-4 framework is national, but the
high-resolution trip log we use here is single-city.  Multi-city
extension is the next step.

\textbf{Zero-inflated target.}  The trip log we used is filtered
to hours with at least one trip per station, which biases the
estimator toward conditional demand rather than full unconditional
demand.  A Hurdle or zero-inflated specification would correct
this.

\textbf{Component T not at national scale.}  The T axis is set to
zero at national application.  The calibration-panel analysis
bounds the information loss to $\sim 8\,\%$.  BD ALTI commune
integration is one engineering task away.

\textbf{Deep-learning scale.}  Three epochs of LSTM training on
CPU are sufficient for the comparison reported here, but a
larger run on GPU with a $96$-hour sequence and a deeper stack is
the proper deep-learning baseline.  Our reading of the recent
tabular-forecasting literature \citep{Grinsztajn2022} is that
this would not change the qualitative ordering, but the gap
would close.

% =========================================================================
\section{Data and code release}
\label{sec:data-release}

The IMD-4 per-commune scores ($n = 34{,}858$) and per-station
scores on the calibration panel ($n = 5{,}442$) are released
under the Open Database Licence at
\url{https://github.com/rohanfosse/imd-national-catalogue}.
The predictive pipeline of Section~\ref{sec:models} -- including
the LightGBM, Ridge, ARIMA, persistence and LSTM model
definitions, the V\'elomagg trip-feature engineering scripts and
the benchmark runner -- is released under MIT at the same URL.
Trip data is derived from the public V\'elomagg dock-based GBFS
feed and the V\'elomagg open-data portal, archived with
SHA-256 hashes for bit-exact reproduction.  Weather features are
from Open-Meteo, archived likewise.

% =========================================================================
\section{Conclusion}
\label{sec:conclusion}

We introduced the IMD-4, a Bayesian composite indicator of
cycling-environment quality computable on every French commune
from public open data, and benchmarked its predictive value
against five baselines on a five-year hourly bike-share trip
panel.  An off-the-shelf LightGBM regressor augmented with the
$13$ IMD station-level features achieves
$R^{2}_{\text{trips}} = 0.311$ on a one-year temporal hold-out,
against $0.042$ for the same regressor without the IMD features
($+0.27$ R$^2$ gap).  The IMD-augmented model matches a
station-fixed-effect baseline within $0.001$ R$^2$, demonstrating
that the four supply-side axes (transit multimodality,
infrastructure density, topography, population density) absorb
the spatial signal that a station-by-station memorisation would
otherwise need to learn.  Classical time-series baselines
(persistence, ARIMA) and a deep-learning baseline (LSTM) are
materially outperformed.

The IMD-4 is computable on all $34{,}858$ French communes, so the
pipeline transfers to any commune with an active dock-based
network.  The natural follow-up is a multi-city benchmark on the
French GBFS portfolio, which we leave for a v1.1 release of the
indicator.

% =========================================================================
\section*{Data and code availability}

The full pipeline -- IMD-4 calibration, national application,
V\'elomagg trip-feature engineering and the six-model benchmark
runner -- is released under MIT at
\url{https://github.com/rohanfosse/imd-national-catalogue}.  The
per-commune IMD-4 scores ($n = 34{,}858$) and per-station scores
on the calibration panel ($n = 5{,}442$) are released as
versioned Parquet files at the same URL under ODbL.  Trained
LightGBM and LSTM model artefacts are released alongside the
benchmark runner.  All raw inputs (V\'elomagg GBFS feed,
data.gouv \emph{lin\'eaire cyclable} and \emph{stations TC},
Open-Meteo weather extracts, INSEE Filosofi and recensement
files) are public and archived with SHA-256 hashes in the same
repository.

\section*{CRediT authorship contribution statement}
\textbf{Rohan Foss\'e :} Conceptualization, Methodology,
Software, Data curation, Formal analysis, Validation,
Visualization, Writing -- original draft, Writing -- review \&
editing.
\textbf{Ga\"el Pallares :} Conceptualization, Methodology,
Supervision, Project administration, Writing -- review \&
editing.

\section*{Declaration of competing interests}
The authors declare that they have no known competing financial
interests or personal relationships that could have appeared to
influence the work reported in this paper.

\section*{Declaration of generative AI and AI-assisted technologies in the writing process}
During the preparation of this work the authors used GitHub
Copilot and Claude (Anthropic) for code completion and editorial
language polishing.  After using these tools the authors reviewed
and edited the content as needed and take full responsibility
for the content of the publication.

\section*{Funding}
This work received no specific grant from any funding agency,
commercial, or not-for-profit entity.  All inputs are public
open data.

\section*{Acknowledgments}
The authors thank the producers of the open data that made this
research possible : data.gouv, MobilityData, OpenStreetMap, INSEE,
Cerema, ONISR, V\'elo \& Territoires, Open-Meteo, and the French
GTFS community ; Montpellier M\'editerran\'ee M\'etropole and the
\emph{Transports de l'agglom\'eration de Montpellier} (TAM) for
the V\'elomagg trip-log open release.

\bibliography{references_demand}

\end{document}
